%LTeX: language=fr-FR
\documentclass{article}
\usepackage{tasks}
\usepackage{amsmath,amssymb,amsfonts}
\usepackage[french]{babel}
\usepackage{enumitem}
\usepackage{currfile}
\usepackage{tikz}
\usetikzlibrary{arrows.meta,calc}
\usepackage{tcolorbox}
\usepackage{fancyhdr}
\usepackage{array}
\usepackage[headsep=3mm,footskip=5mm,includefoot,includehead]{geometry}

\setlength{\parindent}{0cm}

\pagestyle{empty}
\input{\currfiledir ../def.tex}
\def \Document {Activité 1}


\geometry{left=1.5cm, right=1.5cm, top=1cm, bottom=1cm}
\setlength{\headheight}{14.0pt}

\newcolumntype{C}[1]{>{\centering\arraybackslash}m{#1}}

% Déplacement codé en colonne
\newcommand{\dep}[2]{{\renewcommand{\arraystretch}{1}\begin{pmatrix} #1 \\ #2 \end{pmatrix}}}
\newcommand{\depvide}{\begin{pmatrix} \dots \\ \dots \end{pmatrix}}

% Repère et lieux de l'île
\newcommand{\lieu}[3]{\fill (#1) circle (2.5pt); \node[#2, font=\small] at (#1) {#3};}
\newcommand{\croix}[1]{\draw[red!75!black, line width=1.6pt] ($(#1)+(-0.22,-0.22)$) -- ($(#1)+(0.22,0.22)$) ($(#1)+(-0.22,0.22)$) -- ($(#1)+(0.22,-0.22)$);}

% Carte de l'île, avec un dessin en argument
\newcommand{\carte}[1]{%
  \begin{tikzpicture}[x=0.95cm, y=0.95cm]
    \fill[blue!7] (-7.5,-5.5) rectangle (7.5,5.5);
    \foreach \p in {(-6.6,4.9),(6.6,4.9),(-6.6,-5),(6.6,-5),(-3.6,5.2),(3.6,-5.25),(7,0.6),(-7.1,-1.6)} \node[blue!60!black] at \p {$\sim\!\sim$};
    \fill[yellow!22, draw=black, line width=1.2pt] plot[smooth cycle] coordinates
      {(-6.6,-3.4) (-5.2,-4.7) (-1,-4.9) (3,-4.8) (5.6,-3.6) (6.7,-1) (6.2,2.2) (4.6,4.3) (1,4.8) (-2.6,4.5) (-5.6,3.1) (-6.8,0)};
    \draw[gray!45, very thin, step=1] (-7.5,-5.5) grid (7.5,5.5);
    \draw[-{Latex[length=2mm]}] (-7.5,0) -- (7.7,0);
    \draw[-{Latex[length=2mm]}] (0,-5.5) -- (0,5.7);
    \foreach \x in {-7,...,-1,1,2,...,7} \node[below, font=\scriptsize, fill=white, fill opacity=0.6, text opacity=1, inner sep=1pt] at (\x,-0.12) {$\x$};
    \foreach \y in {-5,...,-1,1,2,...,5} \node[left, font=\scriptsize, fill=white, fill opacity=0.6, text opacity=1, inner sep=1pt] at (-0.1,\y) {$\y$};
    \node[below left, font=\scriptsize] at (0,0) {$0$};
    \lieu{-5,-3}{below left}{Ponton $P$}
    \lieu{-4,1}{left}{Rocher $R$}
    \lieu{3,4}{above}{Phare $F$}
    \lieu{-1,2}{above left}{Cocotier $C$}
    \lieu{4,0}{above right}{Grotte $G$}
    \lieu{-4,-1}{left}{Puits $J$}
    \lieu{1,1}{above right}{Cabane $K$}
    \croix{2,-4} \croix{-2,-4} \croix{5,-2} \croix{2,3}
    #1
  \end{tikzpicture}%
}

\begin{document}
\pagestyle{fancy}
\fancyhf{}
\fancyhead[L]{\textbf{\Niveau}}
\fancyhead[C]{\textbf{\ChapitreCourt}}
\fancyhead[R]{\textbf{\Document}}

\large

\begin{center}
  {\Large\bfseries Activité d'introduction : la carte au trésor}
\end{center}

Un pirate a caché son trésor sur une île. La carte est munie d'un repère orthonormé. Sur l'île, il y a quatre croix, mais une seule cache le trésor.

Le pirate code chaque déplacement par une colonne : en haut le déplacement \textbf{horizontal}, en bas le déplacement \textbf{vertical}. Par exemple, $\dep{3}{-2}$ signifie : 3 unités vers la droite, puis 2 unités vers le bas.

\begin{center}
  \carte{}
\end{center}

% =====================================================
\subsection*{Partie A : suivre le message du pirate}

\begin{enumerate}
  \item Le pirate part du ponton $P(-5\,;\,-3)$ et effectue trois déplacements :
  \[
    \text{étape 1 : } \dep{4}{5} \qquad \text{étape 2 : } \dep{5}{-2} \qquad \text{étape 3 : } \dep{-2}{-4}
  \]
  Tracer chaque déplacement par une flèche sur la carte, puis compléter.

  \begin{tasks}[label=\textbullet](3)
    \task Après l'étape 1 : $(\dots\,;\,\dots)$
    \task Après l'étape 2 : $(\dots\,;\,\dots)$
    \task Après l'étape 3 : $(\dots\,;\,\dots)$
  \end{tasks}

  \item Sous quelle croix se trouve le trésor~? \dotfill
\end{enumerate}

\newpage

% =====================================================
\subsection*{Partie B : retrouver le déplacement}

\begin{enumerate}[resume]
  \item Un deuxième pirate va du rocher $R(-4\,;\,1)$ au phare $F(3\,;\,4)$ en un seul déplacement. Lire ce déplacement sur la carte : $\depvide$

  \item Comment peut-on retrouver ce déplacement par un calcul, à partir des coordonnées de $R$ et de $F$~?

  \medskip
  \dotfill

  \item Compléter le tableau par le calcul, sans utiliser la carte.

  \medskip
  \begin{center}
    \renewcommand{\arraystretch}{2.2}
    \begin{tabular}{|C{3.2cm}|C{3.2cm}|C{4.8cm}|}
      \hline
      \textbf{Départ} & \textbf{Arrivée} & \textbf{Déplacement} \\ \hline
      $A(1\,;\,2)$ & $B(4\,;\,7)$ & \\ \hline
      $D(-2\,;\,3)$ & $E(1\,;\,-1)$ & \\ \hline
      $M(3\,;\,-1)$ & $N(-2\,;\,-4)$ & \\ \hline
    \end{tabular}
  \end{center}
\end{enumerate}

% =====================================================
\subsection*{Partie C : le même déplacement}

\begin{enumerate}[resume]
  \item Deux pirates partent l'un du puits $J(-4\,;\,-1)$, l'autre de la cabane $K(1\,;\,1)$. Ils effectuent tous les deux le déplacement $\dep{3}{-2}$.
  Tracer les deux flèches sur la carte et donner les points d'arrivée.

  \medskip
  \dotfill

  \item Que peut-on dire des deux flèches (direction, sens, longueur)~?

  \medskip
  \dotfill
\end{enumerate}

\begin{tcolorbox}[colback=blue!4, colframe=blue!55!black, fonttitle=\bfseries, title=Vocabulaire]
  Un déplacement est représenté par un \textbf{vecteur}. Ses deux nombres, écrits en colonne, sont les \textbf{coordonnées} du vecteur.

  Par exemple, le vecteur $\overrightarrow{RF}$, qui va du rocher au phare, a pour coordonnées $\dep{7}{3}$.

  Deux flèches qui représentent le même déplacement représentent le \textbf{même vecteur} : elles ont la même direction, le même sens et la même longueur.
\end{tcolorbox}

\begin{tcolorbox}[colback=orange!5, colframe=orange!70!black, fonttitle=\bfseries, title=Pour aller plus loin]
  Dans la partie A, le pirate effectue l'étape 1, puis l'étape 2. Quel déplacement unique lui aurait permis d'aller directement du ponton à la grotte~? Comment l'obtenir à partir des déplacements des étapes 1 et 2~?

  \medskip
  \dotfill
\end{tcolorbox}

\end{document}
